% Section 2.4: Summary and Transition to Chapter 3

\section{Summary}
\label{sec:ch2_summary}

This chapter has established the astrophysical and instrumental context for the dark matter searches developed in the remainder of this thesis. The gamma-ray sky observed by the \textit{Fermi}-LAT is a superposition of hadronic and leptonic diffuse emission from cosmic-ray interactions in the Milky Way, thousands of individually resolved point sources---predominantly pulsars and blazars---and a faint, nearly isotropic extragalactic background composed of unresolved active galaxies and star-forming systems. Each of these components is relevant to the analyses that follow: the Galactic Diffuse Emission is the dominant foreground whose imperfect modeling introduces systematic uncertainties in every Fermi-LAT measurement; the spectral degeneracy between millisecond pulsars and dark matter annihilation is at the heart of the Galactic Center Excess debate; and the source-count distribution of unresolved extragalactic populations defines the observable that connects resolved catalogs to the diffuse background.

The \textit{Fermi}-LAT itself imposes fundamental limitations on what can be extracted from the data. The energy-dependent Point Spread Function creates a source confusion limit that prevents standard detection methods from resolving faint objects at low energies. The $TS > 25$ detection threshold partitions the sky into resolved sources entering the catalog and an unresolved component that may contain both astrophysical emitters and dark matter signals. The large population of unassociated sources---over a quarter of the 4FGL catalog---lacks multi-wavelength counterparts and constitutes the primary search pool for exotic emitters such as dark matter subhalos.

These challenges---complex, spatially structured backgrounds; spectral degeneracies between dark matter and conventional sources; instrumental limitations that leave a large fraction of the source population unresolved or unclassified---cannot be addressed by traditional signal-over-background analyses alone. They motivate the statistical and machine learning methods introduced in the next chapter: simulation-based inference to recover population properties below the detection threshold, neural posterior estimation to bypass intractable likelihoods, and domain-adaptive classifiers to operate reliably when training simulations do not perfectly match the observed data.
